% GENERATED FILE. Edit canonical lessons, then run npm run generate:books.
% canonical-source-hash: fnv1a64:2573d4b205497167
% canonical-lessons: HI-C272-gali, HI-C272-nukkar, HI-C272-mela, HI-C272-jalus, HI-C272-bhasan

\chapter{A lane, A street corner, A fair, A procession, A speech}
\label{ch:hi-a-lane-a-street-corner-a-fair-a-procession-a-speech}
\hlchaptermodality{272}

\begin{chapteropening}
\textbf{By the end of this chapter:} \emph{I can say a lane, a street corner, a fair, a procession and a speech.}

Complete the last lesson of chapter 272: I can say a lane, a street corner, a fair, a procession and a speech.
\end{chapteropening}

\section[\texorpdfstring{\hi{गली} (galī)}{गली (galī)}]{\hi{गली} (galī) — a lane}

\label{lesson:HI-C272-gali}



\begin{quote}
Before the new one: say the Hindi for a tenant, then the Hindi for a neighbourhood.
\end{quote}

\subsection*{You'll want to know: \hi{गली}}
\textbf{\hi{गली}} — \emph{galī} — \textquotedblleft{}a lane\textquotedblright{}.

\begin{grammarlens}[title={What you've built}]
One more word for this chapter.
\end{grammarlens}

\subsection*{Your turn}
\begin{itemize}
\raggedright
  \item \emph{Say it:} \emph{galī}
  \item \emph{Say it:} \emph{galī}, once more
  \item \emph{Say it:} say \emph{galī}
  \item \emph{Recall:} say the Hindi for a tenant, then the Hindi for a neighbourhood, then say \emph{galī} again
\end{itemize}

\begin{tcolorbox}[breakable,colback=teal!4,colframe=teal!35!black,title={Before you move on}]
What does \hi{गली} mean? (\textquotedblleft{}A lane\textquotedblright{}.) Say it once more.
\end{tcolorbox}
\section[\texorpdfstring{\hi{नुक्कड़} (nukkaṛ)}{नुक्कड़ (nukkaṛ)}]{\hi{नुक्कड़} (nukkaṛ) — a street corner}

\label{lesson:HI-C272-nukkar}



\begin{quote}
Before the new one: say the Hindi for a neighbourhood, then the Hindi for a lane.
\end{quote}

\subsection*{You'll want to know: \hi{नुक्कड़}}
\textbf{\hi{नुक्कड़}} — \emph{nukkaṛ} — \textquotedblleft{}a street corner\textquotedblright{}.

\begin{grammarlens}[title={What you've built}]
One more word for this chapter.
\end{grammarlens}

\subsection*{Your turn}
\begin{itemize}
\raggedright
  \item \emph{Say it:} \emph{nukkaṛ}
  \item \emph{Say it:} \emph{nukkaṛ}, once more
  \item \emph{Say it:} say \emph{nukkaṛ}
  \item \emph{Recall:} say the Hindi for a neighbourhood, then the Hindi for a lane, then say \emph{nukkaṛ} again
\end{itemize}

\begin{tcolorbox}[breakable,colback=teal!4,colframe=teal!35!black,title={Before you move on}]
What does \hi{नुक्कड़} mean? (\textquotedblleft{}A street corner\textquotedblright{}.) Say it once more.
\end{tcolorbox}
\section[\texorpdfstring{\hi{मेला} (melā)}{मेला (melā)}]{\hi{मेला} (melā) — a fair}

\label{lesson:HI-C272-mela}



\begin{quote}
Before the new one: say the Hindi for a lane, then the Hindi for a street corner.
\end{quote}

\subsection*{You'll want to know: \hi{मेला}}
\textbf{\hi{मेला}} — \emph{melā} — \textquotedblleft{}a fair\textquotedblright{}.

\begin{grammarlens}[title={What you've built}]
One more word for this chapter.
\end{grammarlens}

\subsection*{Your turn}
\begin{itemize}
\raggedright
  \item \emph{Say it:} \emph{melā}
  \item \emph{Say it:} \emph{melā}, once more
  \item \emph{Say it:} say \emph{melā}
  \item \emph{Recall:} say the Hindi for a lane, then the Hindi for a street corner, then say \emph{melā} again
\end{itemize}

\begin{tcolorbox}[breakable,colback=teal!4,colframe=teal!35!black,title={Before you move on}]
What does \hi{मेला} mean? (\textquotedblleft{}A fair\textquotedblright{}.) Say it once more.
\end{tcolorbox}
\section[\texorpdfstring{\hi{जुलूस} (jalūs)}{जुलूस (jalūs)}]{\hi{जुलूस} (jalūs) — a procession}

\label{lesson:HI-C272-jalus}



\begin{quote}
Before the new one: say the Hindi for a street corner, then the Hindi for a fair.
\end{quote}

\subsection*{You'll want to know: \hi{जुलूस}}
\textbf{\hi{जुलूस}} — \emph{jalūs} — \textquotedblleft{}a procession\textquotedblright{}.

\begin{grammarlens}[title={What you've built}]
One more word for this chapter.
\end{grammarlens}

\subsection*{Your turn}
\begin{itemize}
\raggedright
  \item \emph{Say it:} \emph{jalūs}
  \item \emph{Say it:} \emph{jalūs}, once more
  \item \emph{Say it:} say \emph{jalūs}
  \item \emph{Recall:} say the Hindi for a street corner, then the Hindi for a fair, then say \emph{jalūs} again
\end{itemize}

\begin{tcolorbox}[breakable,colback=teal!4,colframe=teal!35!black,title={Before you move on}]
What does \hi{जुलूस} mean? (\textquotedblleft{}A procession\textquotedblright{}.) Say it once more.
\end{tcolorbox}
\section[\texorpdfstring{\hi{भाषण} (bhāṣaṇ)}{भाषण (bhāṣaṇ)}]{\hi{भाषण} (bhāṣaṇ) — a speech}

\label{lesson:HI-C272-bhasan}



\begin{quote}
Before the new one: say the Hindi for a fair, then the Hindi for a procession.
\end{quote}

\subsection*{You'll want to know: \hi{भाषण}}
\textbf{\hi{भाषण}} — \emph{bhāṣaṇ} — \textquotedblleft{}a speech\textquotedblright{}.

\begin{grammarlens}[title={What you've built}]
One more word for this chapter.
\end{grammarlens}

\subsection*{Your turn}
\begin{itemize}
\raggedright
  \item \emph{Say it:} \emph{bhāṣaṇ}
  \item \emph{Say it:} \emph{bhāṣaṇ}, once more
  \item \emph{Say it:} say \emph{bhāṣaṇ}
  \item \emph{Recall:} say the Hindi for a fair, then the Hindi for a procession, then say \emph{bhāṣaṇ} again
\end{itemize}

\begin{tcolorbox}[breakable,colback=teal!4,colframe=teal!35!black,title={Before you move on}]
What does \hi{भाषण} mean? (\textquotedblleft{}A speech\textquotedblright{}.) Say it once more.
\end{tcolorbox}
